\section{Inclusion-exclusion principle}
\begin{theorem}[Exclusion-inclusion]\label{excinc}
Let $A_{1},\dots,A_{n}$ be finite sets. Then the following holds:
$$
\left|\bigcup _{i=1}^{n}A_{i}\right|=\sum_{1\le i\le n} |A_{i}|-\sum _{1\le i <j\le n}|A_{i}\cap A_{j}|+\sum _{1\le i<j<k\le n}|A_{i}\cap A_{j}\cap A_{k}|-\dots+(-1)^{n-1}|A_{1}\cap\dots \cap A_{n}|  
$$
\end{theorem}
\begin{lemma}
Let $B_{1},\dots,B_m$ be finite sets and $w_{1},\dots,w_{m}\in \mathbf{R}$ be some real weights. Then
$$
\sum_{i=1}^{m} w_{i}|B_{i}|=\sum_{i=1}^{m} \sum _{b\in B_{i}}w_{i}=\sum _{b\in \bigcup _{i}B_{i}}\sum _{i:b\in B_{i}}w_{i}
$$
\end{lemma}
\begin{proof}[Proof of the theorem \ref{excinc}]
Our goal is to show that every element of $\bigcup _{i}A_{i}$ is counted in the inclusion-exlcusion formula exactly once.

Suppose that $a\in A_{1}\cup A_{2}\cup\dots \cup A_{n}$ is such that it belongs to precisely $k$ different sets $A_{j},j=1,\dots,n$. This element is counted 
$$
\binom{k}{1}-\binom{k}{2}+\binom{k}{3}+\dots+(-1)^{n-1}\binom{k}{n}=\binom{k}{1}-\binom{k}{2}+\binom{k}{3}+\dots+(-1)^{k-1}\binom{k}{k}=1
$$
time.
\end{proof}
\begin{exercise}[Old exam question]
Consider a $3\times3$ grid with each square colored either blue or red. Compute the total number of colorings which do not contain a $2\times 2$ red square.
\end{exercise}
\begin{proof}[Solution]
We denote $S_{1},S_{2},S_{3},S_{4}$ the sets of colorings that contain a red square. For $j=1,\dots,4$, let $A_{j}$ be the set of colorings that contain the square $S_{j}$. Therefore we can find the number of colorings without square by substracting from the total possible colorings the number of colorings with some squares which gives
$$
\text{\# of colorings without squares }=2^{9}-\bigcup _{j=1}^{4}A_{j}
$$
we then conclude by using \ref{excinc}.
\end{proof}
\section{Applications of the formula}
\begin{question}
A hat-check girl completely loses track of which of $n$ hats belong to which owners, and
hands them back at random to their $n$ owners as the latter leave. What is the
probability $p_{n}$ that nobody receives their own hat back?

Reformulation: Find the number of permutations of the set $[n]$ without fixed points.
\end{question}
\begin{answer}
Let $A$ be the set of all permutations and $A_{i}$ be the set of permutations of the set $[n]$ for which $i$ is a fixed point. The number of permutations with no fixed points is
$$
|A|-\left|\bigcup _{i=1}^{n}A_{i}\right|.
$$
This way $A_{i}\cap A_{j}$ is the set of permutations for which $i$ and $j$ are fixed points and so on:

\begin{align*}
|A| & =n! \\
|A_{i}| & =(n-1)! \\
|A_{j}\cap A_{i}| & =(n-2)! \\
 & \dots
\end{align*}

This gives 

\begin{align*}
|A|-\left|\bigcup _{i=1}^{n}A_{i}\right| & =n!-\binom{n}{1}(n-1)!+\binom{n}{2}(n-2)! \\
 & =n!-\frac{n!(n-1)!}{1!(n-1)!}+\frac{n!(n-2)!}{2!(n-2)!} \\
 & =n!\left( \frac{1}{0!}-\frac{1}{1!}+\frac{1}{2!}-\dots+(-1)^{n} \frac{1}{n!} \right) \\
 & \simeq n!(\exp(-1)).
\end{align*}

Thus we see that the probability $p_{n}$ that nobody receives their own hat back is
$$
p_{n}=\frac{1}{0!}-\frac{1}{1!}+\frac{1}{2!}-\dots+(-1)^{n} \frac{1}{n!}
$$
which as $n$ approaches infinity converges to $\frac{1}{e}$.
\end{answer}

\begin{definition}
In number theory, Euler's totient function $\phi(n)$ counts the positive integers up to a given integer $n$ that are relatively prime to $n$.
\end{definition}
\begin{example}
Among the number $\{ 1,2,3,4,5,6 \}$ only 1 and 5 are coprimes to 6 therefore $\phi(6)=2$.
\end{example}
\begin{example}
If $p$ is a prime number then $\phi(p)=p-1$ and $\phi(p^{k})=p^{k}-p^{k-1}$. 
\end{example}

\begin{proposition}
Suppose that a number $n$ has the prime factorisation $n=p^{k_{1}}_{1}\dots p_{m}^{k_{m}}$ then:
$$
\phi(n)=n\prod_{i=1}^{m} \left( 1-\frac{1}{p_{i}} \right).
$$
\end{proposition}
\begin{proof}
\begin{itemize}
\item Let $A$ be the set of all numbers in $[n]$ not coprime with $n$. 
\item Let $A_{i}$ be the set of all numbers in $[n]$ divisible by $p_{i}$.
\end{itemize}
Then $A=\bigcup _{i=1}^{m}A_{i}$ and $|A_{i}|=\frac{n}{p_{i}}$, $|A_{i}\cap A_{j}|=\frac{n}{p_{i}p_{j}}$ and so on. By \ref{excinc} we find:

\begin{align*}
\phi(n) & =n-|A| \\
 & =n-\sum _{1\le i \le n} \frac{n}{p_{i}}+\sum _{1\le i<j\le m} \frac{n}{p_{i}p_{j}}-\sum _{1\le i<j<k\le m} \frac{n}{p_{i}p_{j}p_{k}}+\dots=n\prod_{i=1}^{m} \left( 1-\frac{1}{p_{i}} \right)
\end{align*}

\end{proof}

\subsection{Bonferroni's inequalities}
\begin{theorem}
Let $A_{1},\dots,A_{n}$ be finite sets. Let $k\in \{ 1,\dots,n \}$ be an odd integer. Then 
$$
\left|\bigcup _{i=1}^{n}A_{i}\right|\le \sum_{s=1}^{k} \sum _{J\subset \{ 1,\dots,n \}}(-1)^{s-1}\cdot\left|\bigcap _{i\in J}A_{j}\right|
$$
Let $\ell \in \{ 1,\dots,n \}$ be an even integer. Then
$$\left|\bigcup _{i=1}^{n}A_{i}\right|\ge \sum _{s=1}^{\ell}\sum _{J\subset \{ 1,\dots,n \}} (-1)^{s-1}\cdot\left|\bigcap j\in J A_{j}\right|$$
\end{theorem}
\begin{proof}
By induction on $n$ we verify the cases $n=1$ and $n=2$ because $|A_{1}|=|A_{1}|$ and $|A_{2}\cup A_{2}|\le |A_{1}|+|A_{2}|$. Therefore we assume the property correct for $n-1$ sets.

Let $k$ be an odd/even integer i $\{ 1,\dots,n-1 \}$:

\begin{align*}
\sum_{s=1}^{k} \sum _{J\subset \{ 1,\dots,n \}}(-1)^{s-1}\left\lvert \bigcup _{j\in J}A_{j}  \right\rvert  & = \\
 & =\sum_{s=1}^{k} \left( \sum _{J\subset \{ 1,\dots,n-1 \}}(-1)^{s-1}\left\lvert  \bigcap _{j\in J}A_{j}  \right\rvert+\sum _{J}(-1)^{s-1}\left\lvert  \bigcap _{j\in J}(A_{j}\cap A_{n})  \right\rvert   \right) \\
 & =\sum_{s=1}^{k} \sum _{J}(-1)^{s-1} \left\lvert  \bigcap _{j\in J}A_{j}  \right\rvert +|A_{n}|-\sum_{s=1}^{k-1} \sum _{J}(-1)^{s}\left\lvert  \bigcap _{j\in J}(A_{j}\cap A_{n})  \right\rvert 
\end{align*}


By the induction hypothesis :
$$
\sum_{s=1}^{k} \sum _{J}(-1)^{s-1}\left\lvert  \bigcap _{j\in J}A_{j}  \right\rvert \le/\ge \left\lvert  \bigcup _{j=1}^{n-1}A_{j}  \right\rvert 
$$
and
$$
\sum_{s=1}^{k-1} \sum _{J}(-1)^{s}\left\lvert  \bigcap _{j\in J}(A_{j}\cap A_{n})  \right\rvert \ge/\le \left\lvert  \bigcup _{j=1}^{n-1}(A_{j}\cap A_{n})  \right\rvert 
$$
Therefore we arrive at
$$
\sum_{s=1}^{k} \sum _{J}(-1)^{s-1}\left\lvert  \bigcap _{j\in J}A_{j}  \right\rvert=\sum_{s=1}^{k} \sum _{J}(-1)^{s-1}\left\lvert  \bigcap _{j\in J}A_{j}  \right\rvert +|A_{n}|-\sum_{s=1}^{k-1} \sum _{J}(-1)^{s}\left\lvert  \bigcap _{j\in J}(A_{j}\cap A_{n})  \right\rvert 
$$
We use the induction hypothesis
$$
\le /\ge \left\lvert  \bigcup _{j=1}^{n-1}A_{j}  \right\rvert +|A_{n}|-\left\lvert  \left( \bigcup _{j=1}^{n-1} \right)\cap A_{n}  \right\rvert .
$$
By \ref{excinc} this number equals:
$$
\left\lvert  \left( \bigcup _{j=1}^{n} \right) \cup A_{n} \right\rvert=\left\lvert  \bigcup _{j=1}^{n}A_{j}  \right\rvert
$$
\end{proof}
